\documentclass[bwprint]{gmcmthesis}
\usepackage{amsmath}
\usepackage{pdfpages}
\numberwithin{figure}{section}
\renewcommand{\thefigure}{\arabic{section}-\arabic{figure}} 
% \documentclass[withoutpreface,bwprint]{cumcmthesis}
% 去掉封面与编号页

\title{水系电解液配方的综合评价、性能预测与实验候选设计}
\baominghao{No.00000000}
\schoolname{匿名学校}
\membera{匿名}
\memberb{匿名}
\memberc{匿名}
\begin{document}
\maketitle
\begin{abstract}
本文针对水系电解液配方的综合评价、性能预测与实验候选设计问题，综合考虑电导率、pH 适宜性、短时电化学稳定性 proxy 以及局部可信域约束，建立了综合评价—树集成预测—可信域分析联合模型，并在 251 条实验记录上完成了从评价到推荐的整体求解。
针对问题一、问题二、问题三，本文分别建立了门槛约束-CRITIC-TOPSIS 综合评价模型、配方结构增强的目标自适应树集成模型和关键组分交互解释模型。其中，问题三显式复用问题二的预测结果与重要性排序，进一步解释高分配方形成机制。
针对问题一，综合评分与电导率排序的 Spearman 相关仅为 0.4154，前 10 名重合仅为 0/10，说明单独使用电导率不足以定义好配方。针对问题二，正式主路线对 conductivity、pH 和 PI 的 OOF MAE 分别为 7.8566、0.1740 和 0.03467，且 PI 直接头与重构头的相关系数达到 0.9662。针对问题三，离子强度 proxy 与加权密度的协同作用进一步解释了高分区形成原因。
在模型检验方面，本文引入结构簇 holdout 与可信域分层。结果表明，low trust 区的 PI 误差约为 high trust 区的 2.57 倍，因此模型更适合在已见邻域内支持开发型决策，而对低可信区域应以探索和风险标注为主。
在此基础上，本文设计了开发型与探索型结合的下一轮实验候选方案，top5 与 top10 的平均预测综合性能分别较随机基线提升 0.0170 和 0.0181。本文的创新点在于：将综合评价、组合预测与可信域分析串联为闭环；把候选推荐与局部扰动稳健性联立审查；并将摘要级硬数字严格限制在 q1、q2、q4、q5 中允许进入摘要的追溯条目内。所建模型可推广至电解液筛选、配方优化及其他小样本实验设计场景。

\keywords{水系电解液\quad 综合评价\quad 树集成预测\quad 可信域\quad 主动实验设计}
\end{abstract}

\tableofcontents

\section{问题重述}

本文围绕水系电解液公开实验数据，依次回答六个相互依赖的子问题。第一阶段关注“如何定义好配方”“如何由配方预测性能”以及“哪些组分和交互作用真正主导结果”；第二阶段进一步追问模型在何处可信、下一轮实验应如何设计，以及高分候选是否处于稳定高性能盆地之中。\par

其中，问题一构造同时覆盖电导率、pH 适宜性和短时稳定性 proxy 的综合性能指标；问题二建立由配方组成直接映射到 conductivity、pH、W\_1、R\_W 与 PI 的预测模型；问题三解释关键组分和交互作用；问题四刻画模型可信域与适用范围；问题五给出下一轮实验的开发型与探索型候选；问题六进一步审查候选在局部扰动下的稳健性。因此，全文并不是孤立回答六个小问，而是在同一数据集上逐步推进实验决策。\par

\section{问题分析}

六个子问题之间不是并列关系，而是严格的链式依赖。问题一先给出统一的综合性能口径；问题二再把配方组成映射到 conductivity、pH、W\_1、R\_W 与 PI；问题三利用问题二的已验证模型解释关键驱动因子和非线性交互；问题四则把误差结构、局部密度与稀有模式合并为可信域分层；在此基础上，问题五据此平衡开发型和探索型实验候选，问题六进一步检验这些候选是在稳定高分区域内，还是只是对扰动敏感的局部尖峰。因此，全文的分析主线不是一次性堆叠多个模型，而是用前一问的证据约束后一问的可用空间和结论强度。\par

\section{模型假设}

\begin{enumerate}
\item 现有 $fast_assessment$ 曲线及其派生量只用于构造短时稳定性 proxy，不将其外推为长期循环寿命。
\item 默认把 pH 工作区间设为 \texttt{[6.5, 8.5]}；区间内记满分，区间外按偏离距离线性递减。
\item 评价指标只基于现有 251 条实验记录构造，因此它代表“当前数据支持下的综合性能口径”，而非普适真值。
\item \texttt{PI}、$W_1$、$R_W$ 的定义完全沿用 \texttt{q1}，因此它们仍然是建立在短时电化学曲线上的 proxy 标签，而不是长期寿命真值。
\item 配方中未出现的组分被记为 \texttt{0.0 structural zero}，代表结构性缺失，而不是随机漏测。
\item 通过 \texttt{volume * density * molality} 构造的 proxy 特征只用于近似投料强度和离子强度，不被解释为严格热力学状态量。
\end{enumerate}

\section{符号说明}

\begin{table}[htbp]
\centering
\small
\begin{tabularx}{\textwidth}{>{\raggedright\arraybackslash}X>{\raggedright\arraybackslash}X}
\toprule
符号 & 含义 \\
\midrule
\texttt{kappa} & 电导率 \\
$S_pH$ & pH 适宜性得分 \\
$W_T = V_c^T - V_a^T$ & Tafel 稳定窗口 \\
$W_1 = V_c^1 - V_a^1$ & $1 mA/cm^2$ 条件下的实用稳定窗口 \\
$R_W = W_1 / W_T$ & 稳定窗口保留率 \\
\texttt{PI} & 综合性能分数 \\
$x_i$ & 第 \texttt{i} 条配方的特征向量 \\
$ratio_j$ & 第 \texttt{j} 种母液的体积分数 \\
$present_j$ & 第 \texttt{j} 种组分是否出现的结构零指示量 \\
$\hat y_i^{(k)}$ & 第 \texttt{k} 个目标的预测值 \\
\bottomrule
\end{tabularx}
\end{table}

\section{模型建立与求解}

\subsection{问题一：综合性能与稳定性指标构造}

\subsubsection{建模动机}

单独使用电导率只能刻画离子传导能力，却无法区分“高电导但酸性过强、实用稳定窗口偏窄”的配方。题面又明确要求结合 pH 与电化学测试曲线讨论更合理的性能表征，因此本问必须把三个维度统一到同一评价框架中。基于此，本文采用门槛约束-CRITIC-TOPSIS 综合评价模型：先把 pH 转换为工作区间适宜性得分，再利用电化学派生量构造稳定窗口指标，最后用 CRITIC 赋权和 TOPSIS 聚合理想解距离。与仅按电导率排序的 baseline 相比，该模型能够直接回答“什么样的配方可以被认为是好的”，并把结果压缩成后续预测问题可用的单一目标量。\par

\subsubsection{核心公式}

本文先对 pH 定义区间适宜性分数：\par

\[
S_{pH}=\max \left(0,\, 1-\frac{d_{pH}}{1.0}\right),
\quad
d_{pH}=\max(0,\, 6.5-pH,\, pH-8.5).
\]

电化学稳定性由两类指标共同表征：\par

\[
W_T = V_c^T - V_a^T,\qquad
W_1 = V_c^1 - V_a^1,\qquad
R_W = \frac{W_1}{W_T}.
\]

设归一化后的第 \texttt{j} 个指标标准差为 $\sigma_j$，相关系数为 $r_{jk}$，则 CRITIC 信息量为\par

\[
C_j = \sigma_j \sum_k (1-r_{jk}), \qquad
w_j = \frac{C_j}{\sum_j C_j}.
\]

将指标向量 $z_i=[kappa_i,S_{pH,i},W_{1,i},R_{W,i}]$ 做 TOPSIS 聚合，得到样本 \texttt{i} 的综合性能分数\par

\[
PI_i=\frac{D_i^-}{D_i^+ + D_i^-},
\]

其中 $D_i^+$ 与 $D_i^-$ 分别表示样本到正理想解和负理想解的距离。\texttt{PI} 越大，样本的综合性能越优。\par

\subsubsection{求解步骤}

\begin{enumerate}
\item 读取 \texttt{A\_data.json} 的 251 条记录，抽取 \texttt{conductivity}、\texttt{pH}、$derived_quantities$ 和配方组成信息。
\item 由派生量计算 $W_T$、$W_1$ 和 $R_W$，并按 \texttt{[6.5, 8.5]} 生成 $S_pH$。
\item 对四个正向指标构造评价矩阵，使用 CRITIC 计算客观权重。
\item 以 TOPSIS 计算综合性能分数 \texttt{PI}，再与电导率单排序、PCA 载荷加权对照进行比较。
\item 将完整中间结果写入 \texttt{indicator\_table.csv}，将可追溯 hard numbers 写入 \texttt{result.json}。
\end{enumerate}

\subsubsection{结果分析}

针对题目中“什么样的配方可以被认为是好的”这一问题，本文认为单独采用电导率作为评价标准并不充分。电导率只能表征离子传导能力，而题面同时强调了 pH 环境与电化学稳定性的重要性，因此我们构造了由电导率 \texttt{kappa}、pH 适宜性 $S_pH$、实用稳定窗口 $W_1$ 与窗口保留率 $R_W$ 组成的综合评价体系。其中，$S_pH$ 用于刻画配方是否处于近中性的可用工作区间，$W_1$ 与 $R_W$ 则分别反映实用电流密度下的稳定窗口大小及其相对于 Tafel 窗口的保留程度。\par

在赋权与聚合方面，本文采用 CRITIC 赋权与 TOPSIS 排序相结合的门槛约束-CRITIC-TOPSIS 综合评价模型。该模型首先根据各指标的离散程度和冲突程度确定客观权重，再通过样本到正负理想解的距离计算综合性能分数 \texttt{PI}。计算结果表明，四项指标的权重依次为 $W_1 = 0.3306$、$S_pH = 0.2703$、$conductivity = 0.2254$、$R_W = 0.1737$，说明在当前数据中，实用稳定窗口和 pH 适宜性对“好配方”的识别作用并不弱于电导率，反而更能区分那些高电导但不适宜实际应用的样本。\par

进一步将综合排序与电导率单排序比较，可见两者的 Spearman 相关系数仅为 \texttt{0.4154}，且前 10 名样本完全不重合。这说明如果只按电导率从高到低筛选配方，将明显偏向酸性较强的样本。事实上，综合评分前 10 样本的平均 pH 为 \texttt{7.833}，平均实用稳定窗口 $W_1$ 为 \texttt{2.5366}；而电导率前 10 样本的平均 pH 仅为 \texttt{5.879}，平均 $W_1$ 也只有 \texttt{1.7432}。因此，从工程可用性的角度看，电导率高并不等于综合性能优。\par

验证与灵敏度分析进一步表明，该综合指标体系在方法层面是可信的，但其精确排名仍具有条件稳定性。一方面，251 条样本均能稳定生成 \texttt{PI}，且窗口指标整体保持正值；另一方面，少量样本的 $R_W$ 略高于 \texttt{1}，提示该指标应被理解为短时 proxy 而非严格机理量。此外，当 pH 工作区间由 \texttt{[6.5, 8.5]} 放宽到 \texttt{[6.0, 8.0]} 时，Top10 与基准只保留 \texttt{6/10}，榜首样本也会变化。因此，本文建议把“需要构造多指标综合评价体系”作为稳健结论写入论文正文，而把“具体哪一条配方排第一”保留为带条件的结果。\par


\begin{figure}[htbp]
\centering
\includegraphics[width=0.82\textwidth]{figures/figure\_q1\_score\_vs\_conductivity.png}
\caption{综合评分与电导率排序差异}
\end{figure}

\subsection{问题二：配方到性能的预测模型}

\subsubsection{建模动机}

题面要求由配方组成信息直接预测电导率、pH 以及 \texttt{q1} 中构造的综合性能与稳定性相关指标，这意味着模型既要容纳配方组分之间的非线性和协同效应，又要在仅有 251 条记录的小样本条件下保持可审计性。基于 \texttt{q2} Plan 与返工 probe，本文采用配方结构增强的目标自适应树路线：先构造由比例、存在性、浓度/质量 proxy、全局聚合特征和配方结构 dummy 组成的统一特征库，再在 \texttt{RandomForest}、\texttt{ExtraTrees}、\texttt{HistGB} 及其简单融合之间进行 OOF 对比，并按照目标差异决定最终路线。为了满足 \texttt{Q2-W05} 的反泄漏约束，正式 Build 没有使用“看测试折调权”的自由融合，而是使用预先固定的简单融合 $Blend_RF_ET$ 与 $Blend_Tree3$。\par

这一路线的选择理由有三点。第一，\texttt{q1} 已经表明单独使用电导率不足以表征“好配方”，因此 \texttt{PI/W\_1/R\_W} 需要与 \texttt{conductivity/pH} 一起被学习。第二，返工 probe 显示简单三树融合在 \texttt{conductivity}、\texttt{pH} 和 \texttt{PI} 三个关键目标上都优于最佳单树，说明融合收益是真实存在的。第三，树路线不仅能输出预测值，还能生成 OOF 残差、切片误差、局部密度和 permutation importance，为 \texttt{q3} 到 \texttt{q6} 提供直接接口。\par

\subsubsection{核心公式}

对每个目标 \texttt{k}，本文先训练三类树基学习器：\par

\[
f_k^{RF}(x), \qquad f_k^{ET}(x), \qquad f_k^{GB}(x).
\]

对于 \texttt{conductivity}、\texttt{pH}、$W_1$ 和 \texttt{PI}，正式 Build 采用预先固定的简单三树融合：\par

\[
\hat y_i^{(k)}=\frac{1}{3}\left(f_k^{RF}(x_i)+f_k^{ET}(x_i)+f_k^{GB}(x_i)\right),
\qquad
k\in\{\text{conductivity},\text{pH},W_1,PI\}.
\]

对于 $R_W$，由于单树 \texttt{RandomForest} 已优于融合候选，因此保留\par

\[
\hat y_i^{(R_W)} = f_{R_W}^{RF}(x_i).
\]

为了检查综合目标量的自洽性，本文还构造\par

\[
\widehat{PI}^{recon}_i = g(\hat \kappa_i,\hat pH_i,\widehat{W_1}_i,\widehat{R_W}_i),
\]

其中 $g(\cdot)$ 使用 \texttt{q1} 固定的 pH 适宜性评分与 CRITIC-TOPSIS 权重口径。若 $\widehat{PI}$ 与 $\widehat{PI}^{recon}$ 偏差过大，则说明直接头虽然精度较高，但解释链条较弱。\par

\subsubsection{求解步骤}

\begin{enumerate}
\item 读取 \texttt{A\_data.json}，为每条配方构造比例特征、存在性特征、\texttt{mass\_proxy\_*}、\texttt{mol\_proxy\_*}、$ionic_strength_proxy$、$weighted_density$、$pattern_key$ 及其 dummy，共得到 \texttt{73} 个输入特征。
\item 读取 \texttt{q1/results/indicator\_table.csv}，按 \texttt{GUID} 合并 \texttt{conductivity}、\texttt{pH}、$W_1$、$R_W$ 与 \texttt{PI} 五个目标。
\item 在 5 折 OOF 下分别训练 \texttt{ElasticNet}、\texttt{PLSRegression}、\texttt{RandomForest}、\texttt{ExtraTrees}、\texttt{HistGB} 及两个固定融合，并比较 OOF 的 MAE、RMSE、$R^2$ 和 Spearman。
\item 根据目标差异确定最终路线：\texttt{conductivity/pH/W\_1/PI -> Blend\_Tree3}，\texttt{R\_W -> RandomForest}。
\item 基于最终路线生成 \texttt{oof\_predictions.csv}、\texttt{slice\_error\_summary.csv}、\texttt{pi\_consistency\_summary.csv}、\texttt{sensitivity\_probe.csv} 和 permutation importance 结果，再把关键硬数字写入 \texttt{result.json}。
\end{enumerate}

\subsubsection{结果分析}

针对题目中“由配方组成信息预测其性能”的要求，本文构建了一个配方结构增强的目标自适应树路线。具体而言，我们首先从 \texttt{A\_data.json} 中提取母液比例、是否出现、浓度/质量 proxy、离子强度 proxy 和配方结构 dummy，共形成 73 个输入特征；随后分别对 \texttt{conductivity}、\texttt{pH}、$W_1$、$R_W$ 与 \texttt{PI} 建立 OOF 预测模型。考虑到样本规模仅有 251 条，且必须避免融合权重信息泄漏，正式实现没有采用基于测试折回填权重的自由 stacking，而是使用了预先固定的简单树融合 $Blend_Tree3$ 与 $Blend_RF_ET$，再结合目标差异选择最终路线。\par

OOF 结果表明，\texttt{conductivity/pH/W\_1/PI} 四个目标都更适合由 $Blend_Tree3$ 预测，而 $R_W$ 保留单树 \texttt{RandomForest} 即可。对应的 OOF MAE 分别为 \texttt{7.8566}、\texttt{0.1740}、\texttt{0.04431}、\texttt{0.01562} 和 \texttt{0.03467}。相较于 \texttt{ElasticNet}，上述主路线在五个目标上的 MAE 都得到降低，其中对综合目标 \texttt{PI} 的改善最为明显，降幅约为 \texttt{24.18\%}。这说明由配方组成到综合性能的映射确实包含不能被单一线性关系充分解释的非线性与交互效应。\par

为了避免把 \texttt{PI} 仅当作一个黑箱分数，本文进一步构造了由预测的 \texttt{conductivity/pH/W\_1/R\_W} 重构得到的 $PI_recon$。结果显示，direct head 的 overall MAE 更低，但它与 $PI_recon$ 的 Spearman 相关仍达到 \texttt{0.9662}。因此，\texttt{PI} 可以继续作为 \texttt{q3} 到 \texttt{q6} 的统一目标量：若更强调预测精度，可优先使用 direct head；若更强调解释链条，则应同时引用 recon 结果。\par

验证与灵敏度分析进一步表明，模型的可信度具有明显区域差异。一方面，删除 proxy block 会使 \texttt{conductivity} 的 MAE 由 \texttt{7.8566} 升至 \texttt{9.5179}，说明模型对浓度强度相关特征具有真实依赖；另一方面，$low_PI_bottom10$ 和 $rare_pattern$ 切片上的误差显著高于整体均值，\texttt{R\_W > 1} 的异常切片更会使 $R_W$ 头明显失稳。由此可见，本文建立的预测模型更适合在已知配方邻域内进行高分区筛选和比较，而不应被润色成“全空间均匀可信”的万能代理。\par


\begin{figure}[htbp]
\centering
\includegraphics[width=0.82\textwidth]{figures/figure\_q2\_model\_mae\_comparison.png}
\caption{主预测模型误差对比}
\end{figure}

\subsection{问题三：关键组分与交互作用解释}

\subsubsection{结果分析}

基于 \texttt{q2} 已通过验证的组合预测模型，我们进一步追问“为什么某些配方更优”。解释结果表明，导电率与综合性能并不由同一组因素完全主导。对 \texttt{conductivity} 而言，\texttt{离子强度 proxy} 及硫酸盐/高氯酸盐相关比例占据头部；对 \texttt{PI} 而言，\texttt{离子强度 proxy}、锂钠体系平衡和最大非水组分占比共同决定了高分区的形成。 进一步的二维非加性分析显示，\texttt{离子强度 proxy} 与 \texttt{加权密度} 不是彼此独立地影响结果，而是在 \texttt{PI} 上形成了显著协同。换言之，某些组分单独看并不突出，但与特定盐型或投料强度 proxy 组合后会显著改变模型输出。 不过，这些规律并非在所有区域同样强。区域一致性统计只给出了 $stable=6$ 个稳定组合，其余 \texttt{4} 个仍需视作条件规律。因此，\texttt{q3} 更适合作为 \texttt{q4} 的可信域切分依据和 \texttt{q5/q6} 的设计线索，而不是把每条解释都包装成普适机理。\par


\begin{figure}[htbp]
\centering
\includegraphics[width=0.82\textwidth]{figures/figure\_q3\_driver\_bars.png}
\caption{关键驱动因子条形图}
\end{figure}

\subsection{问题四：模型可信度与适用范围分析}

\subsubsection{结果分析}

仅靠随机划分训练集与测试集，不足以评价模型在不同配方区域中的真实可用性。我们首先在组合特征空间中把 251 条样本划分为 4 个结构簇，然后分别统计随机 OOF 和 leave-one-cluster-out 的区域误差。结果显示，最差区域的误差倍率可达到整体平均的 \texttt{1.87} 倍，而最差簇外推的 MAE 也明显高于常规随机 OOF。 在此基础上，我们进一步把样本按局部密度、$uncertainty_hook_PI$ 与 rare pattern 标记划分为 high / medium / low trust 三层。\texttt{PI} 在 high trust 区的 MAE 仅为 \texttt{0.0232}，到了 low trust 区则升至 \texttt{0.0597}。因此，\texttt{q2} 模型更适合在 high trust 邻域内做开发型推荐，而对 low trust 区域应以探索和风险标注为主。\par


\begin{figure}[htbp]
\centering
\includegraphics[width=0.82\textwidth]{figures/figure\_q4\_cluster\_error\_heatmap.png}
\caption{结构簇误差热图}
\end{figure}

\subsection{问题五：下一轮实验候选设计}

\subsubsection{结果分析}

在实验预算有限的前提下，下一轮候选设计不能只做“局部最优微调”，也不能完全随机探索。为此，我们以 \texttt{q2} 的预测结果为收益估计，以 \texttt{q4} 的可信域为风险约束，在已观测样本邻域内生成单步与双步体积转移候选。随后把候选分成开发型与探索型两类：前者优先高 $pred_PI$、高 $W_1$、高 pH 适宜度且位于 high/medium trust 邻域；后者允许进入 medium/low trust 区，但必须保持基本性能门槛。 结果表明，top5 推荐的平均 $pred_PI$ 为 \texttt{0.8204}，比同一候选池上的随机基线高 \texttt{0.0170}；top10 推荐的平均 $pred_PI$ 为 \texttt{0.8219}，也高于随机基线 \texttt{0.0181}。因此，当前推荐方案更像“有证据约束的主动实验设计”，而不是盲目穷举或纯贪心搜索。\par


\begin{figure}[htbp]
\centering
\includegraphics[width=0.82\textwidth]{figures/figure\_q5\_candidate\_frontier.png}
\caption{候选前沿与开发-探索平衡}
\end{figure}

\subsection{问题六：候选配方稳健性建模}

\subsubsection{结果分析}

候选配方是否值得继续开发，不能只看中心点预测分数，还必须看其邻域内的表现是否稳定。我们以 \texttt{q5} 的 10 组候选为中心，对每组配方构造半径 \texttt{0.1} 与 \texttt{0.2} 的局部扰动邻域，并比较邻域平均 \texttt{PI}、最小 \texttt{PI}、标准差以及 pH 可接受率。结果显示，在半径 \texttt{0.1} 下共有 \texttt{0} 个候选可归入 stable high-performance basin，而另有 \texttt{5} 个候选虽然中心分数较高，但对局部扰动更敏感，更像 isolated peak。 因此，最终推荐不应简单地“挑最高分”，而应优先选择落在稳定高分盆地内的候选，把孤立尖峰保留为补充探索对象。这一结论与 \texttt{q4} 的可信域分析是一致的：高中心分数只有在局部邻域也稳定时，才更适合进入开发优先序列。\par


\begin{figure}[htbp]
\centering
\includegraphics[width=0.82\textwidth]{figures/figure\_q6\_robustness\_scatter.png}
\caption{候选稳健性散点图}
\end{figure}

\section{结果汇总与验证}

\subsection{最终结果汇总}

\begin{table}[htbp]
\centering
\small
\begin{tabularx}{\textwidth}{>{\raggedright\arraybackslash}X>{\raggedright\arraybackslash}X>{\raggedright\arraybackslash}X>{\raggedright\arraybackslash}X}
\toprule
问题 & 模型 & 核心结论 & 限制 \\
\midrule
q1 & 门槛约束-CRITIC-TOPSIS 综合评价模型 & 综合性能评价需要同时考虑电导率、pH 适宜性和短时稳定性 proxy，单独使用电导率会遗漏一部分高稳定性或近中性样本。 & 稳定性相关指标仅能表征基于 fast\_assessment 电化学曲线的短时稳定性 proxy，不能外推为长期循环寿命结论。 \\
q2 & 配方结构增强的目标自适应受限加权树集成预测模型 & 配方结构增强的树集成主路线在 \texttt{conductivity}、\texttt{pH}、$W_1$、$R_W$ 和 \texttt{PI} 五个目标上均优于线性基线 \texttt{ElasticNet}。 & \texttt{PI}、$W_1$、$R_W$ 仍然继承自 \texttt{q1} 的短时稳定性 proxy 口径，不能写成长期寿命预测。 \\
q3 & 关键组分与交互作用的分区解释模型 & 导电率与综合性能并不是由同一组因素完全主导；离子强度 proxy、硫酸盐/高氯酸盐比例和锂钠体系平衡共同塑造了结果。 & q3 的解释证据建立在 q2 已验证的预测模型与 permutation importance 上，因此属于数据驱动解释，而不是严格机理定律。 \\
q4 & 结构簇-可信域联合验证模型 & 随机训练/测试切分不足以全面评价模型可用性，因为不同结构区域的误差存在显著异质性。 & q4 的可信域仍然来自有限样本上的结构聚类与 OOF 误差，不代表真实工艺边界。 \\
q5 & 可信域约束下的开发-探索联合候选设计模型 & 下一轮实验不应只盯着当前最高分样本附近继续微调，而应在开发型与探索型候选之间保持平衡。 & q5 给出的候选只是一轮实验优先级建议，尚未经过真实实验反馈验证。 \\
q6 & 候选配方邻域稳健性判别模型 & 推荐候选中既存在稳定高分盆地，也存在对小幅扰动敏感的孤立高分点。 & q6 的稳健性仍然是模型预测意义上的局部稳健性，不等于真实工艺鲁棒性。 \\
\bottomrule
\end{tabularx}
\end{table}

摘要中的定量结果仅引用 \texttt{traceability.md} 中 \texttt{allowed in abstract = yes} 的条目，正文不再逐条展示追溯编号与源字段，以保持论文叙事完整。\par

\subsection{q1 综合性能与稳定性指标构造}

验证结果表明，当前指标体系在数值层面是自洽的：251 条样本全部成功生成 $W_T$、$W_1$、$R_W$、$S_pH$ 和 \texttt{PI}，没有缺失值；$W_T$ 与 $W_1$ 全部为正，$S_pH$ 落在 \texttt{[0,1]}，\texttt{PI} 落在 \texttt{[0.1781, 0.8671]}。同时，边界测试表明当四项指标完全相同时，TOPSIS 会退化为并列分数，而不会伪造排序差异；代表性 5 样本 toy demo 也能稳定地产生与单纯电导率不同的排序。\par

需要保留的验证限制有两点。第一，$R_W$ 存在 \texttt{7} 条略高于 \texttt{1.0} 的样本，其中 \texttt{5} 条高于 \texttt{1.05}，最大值为 \texttt{1.0726}，这提示派生阈值口径并非严格满足 \texttt{W\_1 <= W\_T}。第二，稳定性相关指标的物理含义仍应限定为短时 proxy，不能借此直接推出长期寿命。基于这些事实，本问整体 validation verdict 为 \textbf{PASS}，但物理解释强度和精确排名主张需要保留限制。\par

\subsection{q2 配方到性能的预测模型}

验证结果表明，\texttt{q2} 的实现链路是稳定可用的：251 条样本全部成功生成 OOF 预测，没有出现未处理 NaN、负电导率、越界 pH、负 $W_1$ 或超出 \texttt{[0,1]} 的 \texttt{PI} 预测；toy demo 也在 60 样本下成功复现了树路线相对线性基线的优势。更重要的是，\texttt{PI} 直接头与重构头保持了 \texttt{0.9662} 的高 Spearman 相关，不确定性 hook 与误差也存在中等正相关，这意味着主路线既能提供预测值，也能提供后续可信度线索。\par

不过，验证同时揭示了三个不能隐藏的限制。第一，\texttt{conductivity} 头虽然在整体 MAE 上优于 \texttt{ElasticNet}，但在高 \texttt{PI} Top10 切片上的 MAE 为 \texttt{6.7864}，高于线性基线的 \texttt{4.1390}；因此它不是“所有局部区域都更好”的模型。第二，$R_W$ 头是五个目标中最弱的一头，整体 $R^2 = 0.7006$，且在 \texttt{R\_W > 1} 切片上的误差接近整体的 \texttt{3} 倍。第三，\texttt{PI} 的 Top10 只与真实 Top10 重合 \texttt{7/10}，说明它适合候选筛选，但不能被写成“精确复原了真实榜单”。基于这些事实，本问整体 validation verdict 为 \textbf{PASS}，但局部区域和局部目标的 claim 必须显式降级。\par

\subsection{q4 可信域补充}

仅靠随机划分训练集与测试集，不足以评价模型在不同配方区域中的真实可用性。我们首先在组合特征空间中把 251 条样本划分为 4 个结构簇，然后分别统计随机 OOF 和 leave-one-cluster-out 的区域误差。结果显示，最差区域的误差倍率可达到整体平均的 \texttt{1.87} 倍，而最差簇外推的 MAE 也明显高于常规随机 OOF。 在此基础上，我们进一步把样本按局部密度、$uncertainty_hook_PI$ 与 rare pattern 标记划分为 high / medium / low trust 三层。\texttt{PI} 在 high trust 区的 MAE 仅为 \texttt{0.0232}，到了 low trust 区则升至 \texttt{0.0597}。因此，\texttt{q2} 模型更适合在 high trust 邻域内做开发型推荐，而对 low trust 区域应以探索和风险标注为主。\par

\section{灵敏度分析与模型评价}

\subsection{q1 综合性能与稳定性指标构造}

灵敏度分析表明，\textbf{pH 工作区间是当前模型最主要的敏感源}。当区间由 \texttt{[6.5, 8.5]} 改为 \texttt{[6.0, 8.0]} 时，Top10 与基准只保留 \texttt{6/10}，排序相关系数降到 \texttt{0.6936}，榜首样本也发生变化；当区间改为 \texttt{[7.0, 9.0]} 时，Top10 仍保留 \texttt{9/10}，榜首不变。相较之下，权重混合系数 \texttt{alpha} 在 \texttt{[0.6, 1.0]} 内变化时，Top10 至少保留 \texttt{9/10}，Spearman 相关均高于 \texttt{0.99}，说明排序对赋权方式并不脆弱。\par

若只保留 $W_1$ 而移除 $R_W$，Top1 样本仍保持不变，但 Top10 只保留 \texttt{7/10}。这说明 $R_W$ 更多承担“高分样本细排序”的作用，而不是决定是否进入高分区。综合来看，\textbf{“必须采用多指标综合评价”是稳定结论；“哪一条具体配方排第一”则是条件稳定结论}。\par

\subsection{q2 配方到性能的预测模型}

灵敏度分析表明，\textbf{proxy block 是当前主路线最关键的特征组}。当删除 \texttt{mass\_proxy\_*}、\texttt{mol\_proxy\_*}、$weighted_density$、$ionic_strength_proxy$ 等 proxy 特征后，\texttt{conductivity} 的 MAE 从 \texttt{7.8566} 上升到 \texttt{9.5179}，升幅约 \texttt{21.2\%}；\texttt{PI} 也从 \texttt{0.03467} 升到 \texttt{0.03765}，升幅约 \texttt{8.6\%}。这说明模型确实利用了“配方强度”信息，而不是仅靠比例和结构 dummy 在拟合。\par

相比之下，CV 折叠随机种子从 \texttt{42} 改为 \texttt{7} 时，五个目标的 MAE 只发生小幅波动，没有推翻主路线。$Blend_Tree3$ 也在 \texttt{conductivity}、\texttt{pH} 和 \texttt{PI} 上持续优于最佳单树，因此“简单树融合有稳定收益”这一结论是稳健的。\par

真正的失稳边界来自局部区域而不是折叠随机性。以 \texttt{PI} 为例，$low_PI_bottom10$ 切片上的 MAE 为 \texttt{0.1042}，约为整体的 \texttt{3.01} 倍；稀有模式切片上的 MAE 也达到整体的 \texttt{1.53} 倍。若进一步叠加“低 \texttt{PI} + 稀有模式”这类 stress 场景，模型误差会明显放大。因此，本问最稳健的结论不是“模型全空间都很好”，而是“模型在高分区与已见邻域内更可信，在低分区和稀有模式区应谨慎使用”。\par

\subsection{模型优点}

\begin{itemize}
\item 以问题链为主线，避免了多个模型之间口径不一致的常见问题。
\item 把综合评价、预测精度、可信域分层和候选设计串成闭环，便于直接服务实验决策。
\item 所有 hard numbers 均通过 final 层结果和 traceability 约束，避免正文结论与源数据脱钩。
\end{itemize}

\subsection{模型局限与改进}

\begin{itemize}
\item 稳定性相关指标仅能表征基于 fast\_assessment 电化学曲线的短时稳定性 proxy，不能外推为长期循环寿命结论。
\item \texttt{PI}、$W_1$、$R_W$ 仍然继承自 \texttt{q1} 的短时稳定性 proxy 口径，不能写成长期寿命预测。
\item q3 的解释证据建立在 q2 已验证的预测模型与 permutation importance 上，因此属于数据驱动解释，而不是严格机理定律。
\item q4 的可信域仍然来自有限样本上的结构聚类与 OOF 误差，不代表真实工艺边界。
\item q5 给出的候选只是一轮实验优先级建议，尚未经过真实实验反馈验证。
\item q6 的稳健性仍然是模型预测意义上的局部稳健性，不等于真实工艺鲁棒性。
\item 后续若获得更长时程的稳定性或寿命数据，应重建 \texttt{W\_1 / R\_W / PI} 的标签语义，再重新训练问题二到问题六的链条。
\item Stage 8 的最终论文可继续根据 \texttt{paper.md} 做人工润色，但不应再回退为审查材料拼接稿。
\end{itemize}

\section{参考文献}

\begin{enumerate}
\item Hwang C L, Yoon K. Multiple Attribute Decision Making: Methods and Applications[M]. New York: Springer, 1981.
\item Diakoulaki D, Mavrotas G, Papayannakis L. Determining objective weights in multiple criteria problems: The CRITIC method[J]. Computers and Operations Research, 1995, 22(7): 763-770.
\item Breiman L. Random Forests[J]. Machine Learning, 2001, 45(1): 5-32.
\item Geurts P, Ernst D, Wehenkel L. Extremely randomized trees[J]. Machine Learning, 2006, 63(1): 3-42.
\item Friedman J H. Greedy function approximation: A gradient boosting machine[J]. Annals of Statistics, 2001, 29(5): 1189-1232.
\item 本题公开问题说明与附件数据：\texttt{problem.md}、\texttt{A\_data.json}、\texttt{README.txt}。
\end{enumerate}

\section{附录}

\subsection{附录 A 主要程序说明}

\begin{itemize}
\item \texttt{q1\_build.py}：生成综合性能指标、权重与灵敏度分析。
\item \texttt{q2\_build.py}：训练配方到性能的主预测模型并输出 OOF 结果。
\item \texttt{q3\_build.py}：提取关键驱动因子与交互作用解释。
\item \texttt{q4\_build.py}：执行结构簇验证与可信域分层。
\item \texttt{q5\_build.py}：构造下一轮实验候选池并完成开发-探索推荐。
\item \texttt{q6\_build.py}：对候选进行局部扰动稳健性审查。
\end{itemize}

\subsection{附录 B 结果材料说明}

\begin{itemize}
\item 全部图表素材已经复制到 \texttt{final/source/figures/} 与 \texttt{final/source/tables/}。
\item 详细追溯链保留在 \texttt{traceability.md}，不再直接写入正文。
\end{itemize}
\end{document}
